% Libro 1 · Mundo 1 · Cálculo mental (app: pu-calculo)
\mundo{Cálculo mental}{pu-calculo}{En el examen no hay calculadora, pero tampoco hacen falta cuentas largas: las preguntas
están pensadas para quien ve la estructura de una operación antes de hacerla. Este mundo repasa los números con los que
se trabaja todo el libro (enteros, fracciones, decimales, porcentajes y potencias) y, sobre todo, enseña a calcular
\emph{poco} y \emph{bien}.}{%
\item \textbf{1.1} Enteros, signos y jerarquía de operaciones
\item \textbf{1.2} Fracciones: partir, comparar y operar
\item \textbf{1.3} Decimales, periodos y porcentajes
\item \textbf{1.4} Potencias, raíces y estimación}

% ------------------------------------------------------------------
\leccion{pu-cm-1}{Enteros, signos y jerarquía}{Video «Enteros» y los juegos del globo con lastre, la flecha que se voltea y
las capas de una operación.}

\capitulo{La recta}
Un entero no es solo una cantidad: es un \clave{lugar}. Piensa en un termómetro acostado. El 0 es la frontera; a la derecha
viven los positivos y a la izquierda los negativos.

\xrecta{-8}{8}{\xficha[zfire]{-6}\xficha[zfire]{-1}\xficha{4}}
\pie{$-6 < -1 < 4$: es menor el que está más a la izquierda.}

¿Qué es menor, $-6$ o $-1$? El número 6 es más grande que el 1, pero aquí no se comparan tamaños sino lugares: $-6$ está más a
la izquierda, así que $-6 < -1$. Si $-6$ y $-1$ son temperaturas, $-6$ es más frío.

\turno{Ubica en la recta los números $-7$, $2$, $-3$ y $6$, y escríbelos de menor a mayor.
\xrecta{-8}{8}{}}{$-7 < -3 < 2 < 6$.}

\textbf{El opuesto es un reflejo.} Imagina un espejo en el 0. El reflejo de $5$ es $-5$ y el reflejo de $-5$ es $5$. El signo
menos delante de un número \emph{pide} su reflejo: por eso $-(-5) = 5$, y por eso un número más su opuesto da 0.

\xrecta{-8}{8}{\draw[zlime,line width=1.6pt] (0,-.5)--(0,.95);\xficha{5}\xficha[zfire]{-5}\xsalto[zmute]{5}{-5}{reflejo}}

\textbf{La distancia cuenta pasos.} De $-3$ a $5$ se camina 3 pasos hasta el 0 y 5 más: son 8. Se escribe con barras,
$|5 - (-3)| = 8$, y nunca es negativa, porque los pasos se cuentan, no se deben. Como se puede caminar hacia los dos
lados, a una misma distancia de un número hay \clave{dos} lugares.

\turno{¿Qué lugares están a 6 pasos de $-1$? Márcalos. \xrecta{-8}{8}{\xficha{-1}}}{$-7$ y $5$: $-1 - 6$ y $-1 + 6$.}

\capitulo{Sumar y restar}
\textbf{Sumar es caminar.} El signo dice hacia dónde y el número, cuántos pasos: sumar un positivo lleva a la derecha y sumar
un negativo, a la izquierda.

\xrecta{-8}{8}{\xficha{-5}\xficha{3}\xsalto{-5}{3}{$+8$}}
\cuenta{-5 + 8 = 3}

\textbf{Restar es caminar al revés.} Sumar $-5$ llevaría 5 pasos a la izquierda; restar $-5$ es hacer lo contrario: 5 pasos a
la derecha. Piénsalo como dinero: si te perdonan una deuda de 5, tu saldo \emph{sube} 5.

\xrecta{-8}{8}{\xficha{2}\xficha{7}\xsalto{2}{7}{$-\,(-5)$}}
\cuenta{2 - (-5) = 2 + 5 = 7}

\turno{Dibuja el salto y escribe dónde caes: $\;-1 - 6\;$ y luego $\;-3 - (-9)$.
\xrecta{-8}{8}{}}{$-1 - 6 = -7$ (seis pasos a la izquierda) y $-3 - (-9) = 6$ (nueve a la derecha).}

\capitulo{El menos da la vuelta}
\textbf{Multiplicar es repetir saltos iguales}, como una rana que siempre salta lo mismo. $4 \cdot (-2)$ son cuatro saltos de
$-2$: llegan a $-8$. Es como deber 2 a cuatro amigos.

\xrecta{-9}{9}{\foreach \a/\b in {0/-2,-2/-4,-4/-6,-6/-8}{\draw[zfire,line width=1.1pt,-{Stealth[length=2mm]}] (\a,.12) to[bend right=55] (\b,.12);}
\foreach \a/\b in {0/2,2/4,4/6,6/8}{\draw[zolive,line width=1.1pt,dashed,-{Stealth[length=2mm]}] (\a,.12) to[bend left=55] (\b,.12);}}
\pie{Naranja: $4 \cdot (-2) = -8$. Al cambiar el 4 por $-4$, toda la cadena da la vuelta en el 0: $(-4)(-2) = 8$.}

Un factor negativo \clave{da la vuelta} a la cadena. Dos vueltas te dejan como estabas: por eso menos por menos da más. Con
muchos factores no importa \emph{cuál} es negativo, sino \emph{cuántos}: par, positivo; impar, negativo.

\turno{Sin multiplicar, di el signo de $(-3)(-1)(2)(-5)$ y el de $(-4)(6)(-2)(1)$.}{El primero, negativo (tres
vueltas); el segundo, positivo (dos vueltas).}

\capitulo{Las capas}
En una cuenta larga el orden importa. Se resuelve por capas, de adentro hacia afuera, como pelar una cebolla.

\begin{center}\begin{tikzpicture}[x=1mm,y=1mm,font=\sffamily\small]
\foreach \y/\e/\r in {0/{$7 - 2 \cdot (\,1 - 5\,)^2$}/{paréntesis: $1 - 5 = -4$},-10/{$7 - 2 \cdot (-4)^2$}/{potencia: $(-4)^2 = 16$},
  -20/{$7 - 2 \cdot 16$}/{producto: $2 \cdot 16 = 32$},-30/{$7 - 32 = -25$}/{resta}}{%
  \node[fill=zsoft,rounded corners=1.2mm,minimum width=58mm,minimum height=8mm] at (0,\y) {\e};
  \node[anchor=west,text=zmute] at (34,\y) {\r};}
\end{tikzpicture}\end{center}

\turno{Escribe una capa por renglón: $\;10 + 3 \cdot (2 - 6) \div 2$.\renglones{3}}{$10 + 3 \cdot (-4) \div 2 = 10 - 12 \div 2 = 10 - 6 = 4$.}

\begin{regla}
\begin{itemize}
  \item Restar es sumar el opuesto: $a - b = a + (-b)$; en particular $a - (-b) = a + b$.
  \item Distancia entre $a$ y $b$: $|a - b|$. Hay dos números a distancia $d$ de $a$: $a - d$ y $a + d$.
  \item En productos y cocientes: si hay una cantidad \textbf{par} de factores negativos, el resultado es positivo; si es
  \textbf{impar}, es negativo.
  \item Orden de las operaciones (de adentro hacia afuera): \textbf{paréntesis} $\to$ \textbf{potencias y raíces}
  $\to$ \textbf{multiplicaciones y divisiones} (de izquierda a derecha) $\to$ \textbf{sumas y restas} (de izquierda a
  derecha).
\end{itemize}
\end{regla}

\subsection{Ejemplos resueltos}
\begin{ejemplo}[Contar signos]
Calcula $\dfrac{(-2)(-6)(-5)}{(-3)(-4)}$.
\solucion Primero el signo: en total hay cinco signos menos (tres arriba y dos abajo), un número impar, así que el
resultado es negativo. Luego el tamaño, simplificando antes de multiplicar: $\dfrac{2 \cdot 6 \cdot 5}{3 \cdot 4} =
\dfrac{60}{12} = 5$. El resultado es $-5$.
\end{ejemplo}

\begin{ejemplo}[Capa por capa]
Calcula $20 - 3\,(4 - 7)^2 \div 9 + (-2)^3$.
\solucion
\paso{Paréntesis.} $4 - 7 = -3$.
\paso{Potencias.} $(-3)^2 = 9$ y $(-2)^3 = -8$. Queda $20 - 3 \cdot 9 \div 9 + (-8)$.
\paso{× y ÷ de izquierda a derecha.} $3 \cdot 9 = 27$ y $27 \div 9 = 3$. Queda $20 - 3 - 8$.
\paso{+ y −.} $20 - 3 - 8 = 9$.
\end{ejemplo}

\begin{ejemplo}[Los dos lugares a la misma distancia]
Un número está a distancia 7 de $-2$ y es negativo. ¿Cuál es?
\solucion A distancia 7 de $-2$ están $-2 - 7 = -9$ y $-2 + 7 = 5$. El negativo es $-9$.
\end{ejemplo}

\begin{ejemplo}[Pisos y sótanos]
Un ascensor parte del sótano 3 (piso $-3$), sube 9 pisos, baja 4 y vuelve a subir 6. Si cada piso mide 3 metros, ¿a
cuántos metros del punto de partida terminó?
\solucion El recorrido neto es $+9 - 4 + 6 = +11$ pisos: termina en el piso $-3 + 11 = 8$. La distancia al punto de
partida es $11 \cdot 3 = 33$ metros. Observa que no hizo falta saber en qué piso terminó: bastaba el desplazamiento neto.
\end{ejemplo}

\begin{trampa}
$-3^2$ \textbf{no} es lo mismo que $(-3)^2$. La potencia se hace antes que el signo menos: $-3^2 = -(3^2) = -9$, mientras
que $(-3)^2 = 9$. Las opciones del examen suelen traer las dos.
\end{trampa}

\begin{atajo}
Separa siempre \emph{el signo} del \emph{tamaño}. Cuenta los menos para el signo y luego simplifica los números en valor
absoluto. Así nunca multiplicas números grandes con signos mezclados.
\end{atajo}

\begin{profundiza}
La jerarquía no es un capricho: es una convención para que una expresión escrita tenga un solo valor. Lo mismo pasa con
«de izquierda a derecha»: $24 \div 4 \cdot 2$ vale $12$, no $3$, porque la división va primero por estar a la izquierda.
\end{profundiza}

\begin{reto}
Un congelador pasa de $-18\,^{\circ}\mathrm{C}$ a $4\,^{\circ}\mathrm{C}$ en 11 minutos, subiendo siempre al mismo
ritmo. ¿Qué temperatura marcaba a los 3 minutos?\par
\opciones{$-12\,^{\circ}\mathrm{C}$}{$-14\,^{\circ}\mathrm{C}$}{$-6\,^{\circ}\mathrm{C}$}{$-15\,^{\circ}\mathrm{C}$}
\end{reto}

\begin{retoresuelto}
De $-18$ a $4$ hay $4 - (-18) = 22$ grados, que se suben en 11 minutos: $2$ grados por minuto. A los 3 minutos subió
$6$ grados: $-18 + 6 = -12\,^{\circ}\mathrm{C}$. \textbf{Opción A.} La D sale de suponer un grado por minuto ($-18 + 3$) y la B de sumar al inicio los 4 grados finales ($-18 + 4$).
\end{retoresuelto}

\armar{
\arma{El viaje}{Vas de $-9$ a $4$ con \emph{dos} movimientos: uno debe ser restar un negativo. Escríbelos y dibuja los saltos.}{Hay muchas:
por ejemplo $-9 - (-15) - 2 = 4$. Los dos movimientos deben subir 13 en total. Error frecuente: creer que $-\,(-15)$ baja.}
\arma{Los signos}{Encuentra dos enteros cuyo producto sea positivo y cuya suma sea negativa. ¿Cómo son siempre?}{Los dos negativos,
como $-2$ y $-7$. Error frecuente: usar para la suma la regla de signos del producto.}
\arma{El paréntesis justo}{Pon un solo paréntesis en $\;6 - 2 \cdot 5 - 8\;$ para que el resultado sea $12$.}{$6 - 2 \cdot (5 - 8) =
6 - 2 \cdot (-3) = 6 + 6 = 12$. Error frecuente: dejar $6 - 6$ al multiplicar los signos.}
}

\taller
\piso{Calienta}{operaciones directas}
\begin{ejercicios}[2]
\item $-7 + 12 - 9$ \sol{$-4$.}
\item $5 - (-8) - 3$ \sol{$10$.}
\item $(-4)(-5)(-1)$ \sol{$-20$ (tres signos menos: impar).}
\item $-2^4 + (-2)^4$ \sol{$-16 + 16 = 0$.}
\item $|{-6} - 9|$ \sol{$15$.}
\item $36 \div (-4) \cdot 3$ \sol{$-27$: de izquierda a derecha, $-9 \cdot 3$.}
\end{ejercicios}
\piso{Entrena}{problemas cortos}
\begin{ejercicios}
\item Calcula $\;8 - 2\,[\,3 - (5 - 9)\,] \div 7$. \sol{$8 - 2 \cdot 7 \div 7 = 8 - 2 = 6$.}
\item Un buzo está a $-14$ m. Sube 9 m, baja 23 m y sube 11 m. ¿A qué profundidad quedó? \sol{$-14 + 9 - 23 + 11 = -17$: a 17 m de profundidad.}
\item ¿Cuántos enteros hay estrictamente entre $-8$ y $5$? \sol{De $-7$ a $4$: $4 - (-7) + 1 = 12$ enteros.}
\item Dos números enteros están a distancia 10 entre sí y su suma es 4. ¿Cuáles son? \sol{Su punto medio es $2$; están a 5 de él: $-3$ y $7$.}
\item Ayer a las 6 a.\,m. hacía $-5\,^{\circ}\mathrm{C}$ y al mediodía $13\,^{\circ}\mathrm{C}$. ¿Cuántos grados subió por hora, si subió parejo? \sol{$18$ grados en 6 horas: 3 por hora.}
\end{ejercicios}
\piso{Como en el examen}{opción múltiple}
\begin{ejercicios}
\item El valor de $\;-(-2)^2 - 2^2 \cdot (-1)^{5}\;$ es
\opciones{$-8$}{$0$}{$8$}{$-4$}
\sol{B. $-(4) - 4 \cdot (-1) = -4 + 4 = 0$.}
\item Tres amigos anotan las ganancias de un juego: Pedro $-7$, Ana $+12$ y Lucía. Si la suma de los tres es $-3$,
¿cuánto le falta a Lucía para llegar a 0?
\opciones{$8$}{$5$}{$-8$}{$2$}
\sol{A. Lucía tiene $-3 - (-7) - 12 = -8$; le faltan 8.}
\item Si $a < 0 < b$, ¿cuál de las siguientes expresiones es \textbf{siempre} negativa?
\opciones{$a + b$}{$b - a$}{$a \cdot b$}{$a^2 \cdot b$}
\sol{C. Negativo por positivo es negativo. $a+b$ depende de los tamaños; $b - a > 0$; $a^2 b > 0$.}
\end{ejercicios}

\bitacora

% ------------------------------------------------------------------
\leccion{pu-cm-2}{Fracciones}{Video «Fracciones» y los juegos de cortar, comparar y repartir lo que queda.}

\capitulo{Partes iguales}
Una fracción aparece cuando una barra se corta en pedazos \clave{iguales}. El número de abajo, el \clave{denominador}, dice en
cuántos se cortó; el de arriba, el \clave{numerador}, cuántos se toman.

\begin{barras}\xbarra{0}{9}{5}{$\tfrac59$}\xbarra{-12}{9}{9}{$\tfrac99 = 1$}\end{barras}

Si un pedazo fuera más grande que otro, «cinco de nueve» no diría cuánto es. Y si se toman todos, vuelve la barra entera.

\turno{Corta la barra y sombrea $\tfrac58$. ¿Qué fracción queda en blanco?
\begin{barras}\xbarra{0}{1}{0}{}\end{barras}}{Queda $\tfrac38$: $\tfrac58 + \tfrac38 = 1$.}

\capitulo{Fracciones equivalentes}
Cortar cada pedazo en dos duplica los pedazos \emph{y} los que se toman: cambia el nombre, no la cantidad. Simplificar es el
camino de vuelta: dividir arriba y abajo entre el mismo número.

\begin{barras}\xbarra{0}{3}{2}{$\tfrac23$}\xbarra{-9}{6}{4}{$\tfrac46$}\xbarra{-18}{12}{8}{$\tfrac{8}{12}$}\end{barras}

\turno{Completa: $\;\dfrac{3}{5} = \dfrac{\rule{7mm}{.4pt}}{20}\;$ y simplifica $\;\dfrac{18}{24}$.}{$\tfrac{12}{20}$ y $\tfrac34$.}

\capitulo{Comparar y sumar}
Pedazos de distinto tamaño no se pueden contar juntos. El truco es volver a cortar hasta que \clave{midan lo mismo}.

\begin{barras}\xbarra{0}{4}{1}{$\tfrac14$}\xbarra[zfire]{-10}{6}{1}{$\tfrac16$}
\xbarra{-24}{12}{3}{$\tfrac{3}{12} + \tfrac{2}{12} = \tfrac{5}{12}$}
\foreach \i in {3,4}{\filldraw[fill=zfire!75,draw=zmute,line width=.5pt] (\i*8,-24) rectangle ++(8,8);}\end{barras}

Con pedazos iguales se cuenta y ya: 3 más 2 son 5 doceavos. Sumar arriba y abajo daría $\tfrac{2}{10}$, que es \emph{menos} que
$\tfrac14$: absurdo. Para comparar, el atajo es el \clave{producto cruzado}: $\tfrac47 > \tfrac59$ porque $4 \cdot 9 = 36$ gana a
$5 \cdot 7 = 35$. Es lo mismo que pasar las dos a pedazos de $\tfrac{1}{63}$.

\turno{¿Cuál es mayor, $\tfrac{5}{8}$ o $\tfrac{7}{11}$? ¿Y cuánto es $\tfrac13 + \tfrac25$?}{$5 \cdot 11 = 55 < 7 \cdot 8 = 56$: es mayor
$\tfrac{7}{11}$. La suma es $\tfrac{5}{15} + \tfrac{6}{15} = \tfrac{11}{15}$.}

\capitulo{La fracción de una cantidad}
Ahora la barra \emph{vale} algo. Para $\tfrac27$ de 35 se reparten 35 en 7 grupos (de 5) y se toman 2: son 10.

\begin{barras}\xbarra{0}{7}{2}{$\tfrac27$ de $35 = 10$}\foreach \i in {1,...,7}{\node[font=\sffamily\footnotesize] at ({(\i-.5)*96/7},4) {5};}\end{barras}

La palabra «de» multiplica, y también vale con otra fracción: $\tfrac23$ \emph{de} $\tfrac34$ es
$\tfrac23 \cdot \tfrac34 = \tfrac12$. Dividir pregunta lo contrario, \emph{cuántas caben}: en 3 caben doce cuartos, así que
$3 \div \tfrac14 = 3 \cdot 4 = 12$.

\turno{En un bus de 40 puestos van ocupados tres octavos. Dibuja la barra y di cuántos puestos quedan libres.
\begin{barras}\xbarra{0}{8}{0}{}\end{barras}}{Cada octavo son 5 puestos: van 15 y quedan 25.}

\begin{regla}
\[\frac ab + \frac cd = \frac{ad + bc}{bd}\qquad \frac ab \cdot \frac cd = \frac{ac}{bd}\qquad
\frac ab \div \frac cd = \frac ab \cdot \frac dc\]
\[\frac ab < \frac cd \iff ad < bc\quad (b, d > 0)\]
\textbf{Fracción de lo que queda:} si se usa $\tfrac{p}{q}$ de algo, queda $1 - \tfrac pq$. Si luego se usa
$\tfrac rs$ \emph{del resto}, queda $\left(1 - \tfrac pq\right)\left(1 - \tfrac rs\right)$ del total.
\end{regla}

\subsection{Ejemplos resueltos}
\begin{ejemplo}[Simplificar antes de multiplicar]
Calcula $\dfrac{14}{15} \cdot \dfrac{25}{21} \cdot \dfrac{9}{10}$.
\solucion No se multiplica nada: se descompone en factores primos y se cancela.
Numerador: $14 \cdot 25 \cdot 9 = 2\cdot7\cdot5\cdot5\cdot3\cdot3$. Denominador:
$15 \cdot 21 \cdot 10 = 3\cdot5\cdot3\cdot7\cdot2\cdot5$. Son exactamente los mismos factores, así que el producto es $1$.
\end{ejemplo}

\begin{ejemplo}[La fracción del resto]
Marta gasta $\tfrac14$ de su sueldo en arriendo y $\tfrac25$ \emph{de lo que le queda} en comida. Si le sobran
\$540\,000, ¿cuál es su sueldo?
\solucion Tras el arriendo queda $\tfrac34$. De eso gasta $\tfrac25$, así que conserva $\tfrac35$ de $\tfrac34$:
$\tfrac35 \cdot \tfrac34 = \tfrac{9}{20}$ del sueldo. Si $\tfrac{9}{20}$ son 540\,000, $\tfrac{1}{20}$ es 60\,000 y el
sueldo es $20 \cdot 60\,000 = 1\,200\,000$.
\end{ejemplo}

\begin{ejemplo}[Ordenar sin decimales]
Ordena de menor a mayor $\tfrac58$, $\tfrac23$ y $\tfrac{7}{12}$.
\solucion Denominador común 24: $\tfrac{15}{24}$, $\tfrac{16}{24}$, $\tfrac{14}{24}$. Orden:
$\tfrac{7}{12} < \tfrac58 < \tfrac23$.
\end{ejemplo}

\begin{ejemplo}[Fracciones de fracciones en un problema]
En un salón, $\tfrac23$ de los estudiantes son mujeres y $\tfrac38$ de las mujeres usan gafas. ¿Qué fracción del salón son
mujeres sin gafas?
\solucion Sin gafas son $\tfrac58$ de las mujeres: $\tfrac58 \cdot \tfrac23 = \tfrac{10}{24} = \tfrac{5}{12}$ del salón.
\end{ejemplo}

\begin{trampa}
En la fracción del resto, el error típico es restar las fracciones del total: «$1 - \tfrac14 - \tfrac25$». La segunda
fracción es de \emph{lo que quedaba}, no del total; por eso se multiplica.
\end{trampa}

\begin{atajo}
Si en un problema «$\tfrac{9}{20}$ del total son 540\,000», no plantees ecuación: divide por el numerador para hallar una
parte ($60\,000$) y multiplica por el denominador. Es la misma cuenta, más limpia.
\end{atajo}

\begin{reto}
Una jarra está llena hasta los $\tfrac35$. Se le echan 2 vasos de agua y queda llena hasta los $\tfrac45$. ¿Cuántos vasos
caben en la jarra completa?
\opciones{8}{10}{12}{15}
\end{reto}

\begin{retoresuelto}
Los 2 vasos llenan la diferencia entre $\tfrac45$ y $\tfrac35$, que es $\tfrac15$ de la jarra. Si un quinto son 2 vasos,
la jarra completa son $5 \cdot 2 = 10$ vasos. \textbf{Opción B.}
\end{retoresuelto}

\armar{
\arma{Entre dos}{Escribe una fracción que esté entre $\tfrac25$ y $\tfrac12$. Muestra con dos barras por qué sirve.}{Por ejemplo
$\tfrac{9}{20}$: $\tfrac25 = \tfrac{8}{20}$ y $\tfrac12 = \tfrac{10}{20}$. Error frecuente: creer que no hay nada entre ellas.}
\arma{La suma justa}{Encuentra dos fracciones con denominadores distintos que sumen $\tfrac{7}{12}$.}{Por ejemplo
$\tfrac13 + \tfrac14$. Error frecuente: sumar arriba y abajo.}
\arma{El total escondido}{Inventa una cantidad de la que $\tfrac35$ sea 27 y explica con una barra cómo la hallaste.}{Tres quintos
son 27, cada quinto es 9 y el total es 45. Error frecuente: calcular $\tfrac35$ de 27.}
}

\taller
\piso{Calienta}{operaciones}
\begin{ejercicios}[2]
\item $\tfrac23 + \tfrac34$ \sol{$\tfrac{17}{12}$.}
\item $\tfrac56 - \tfrac49$ \sol{$\tfrac{15 - 8}{18} = \tfrac{7}{18}$.}
\item $\tfrac{12}{35} \cdot \tfrac{14}{9}$ \sol{$\tfrac{8}{15}$.}
\item $\tfrac34 \div \tfrac{9}{16}$ \sol{$\tfrac34 \cdot \tfrac{16}{9} = \tfrac43$.}
\item $\tfrac35$ de $45$ \sol{$27$.}
\item Simplifica $\tfrac{84}{126}$ \sol{$\tfrac23$ (se divide por 42).}
\end{ejercicios}
\piso{Entrena}{problemas}
\begin{ejercicios}
\item ¿Qué fracción es mayor, $\tfrac{11}{13}$ o $\tfrac{13}{15}$? \sol{$\tfrac{13}{15}$: $11 \cdot 15 = 165 < 169 = 13 \cdot 13$. A cada una le falta $\tfrac{2}{13}$ y $\tfrac{2}{15}$ para el 1, y le falta menos a la segunda.}
\item Un tanque tiene $\tfrac38$ de agua. Se gastan 30 litros y queda en $\tfrac14$. ¿Cuál es su capacidad? \sol{$\tfrac38 - \tfrac14 = \tfrac18$ son 30 L: 240 L.}
\item Juan leyó $\tfrac13$ de un libro el lunes y $\tfrac14$ del resto el martes. Le faltan 150 páginas. ¿Cuántas tiene el libro? \sol{Queda $\tfrac23 \cdot \tfrac34 = \tfrac12$: el libro tiene 300 páginas.}
\item ¿Cuántos trozos de $\tfrac38$ m se cortan de una cinta de $6$ m? \sol{$6 \div \tfrac38 = 16$.}
\end{ejercicios}
\piso{Como en el examen}{opción múltiple}
\begin{ejercicios}
\item El valor de $\left(1 - \tfrac12\right)\left(1 - \tfrac13\right)\left(1 - \tfrac14\right)\cdots\left(1 - \tfrac{1}{20}\right)$ es
\opciones{$\tfrac{1}{20}$}{$\tfrac{19}{20}$}{$\tfrac{1}{19}$}{$\tfrac{1}{2}$}
\sol{A. $\tfrac12 \cdot \tfrac23 \cdot \tfrac34 \cdots \tfrac{19}{20}$: cada numerador se cancela con el denominador anterior.}
\item De un grupo, $\tfrac25$ viajan en bus, $\tfrac13$ en metro y los 12 restantes a pie. ¿Cuántos son?
\opciones{30}{45}{60}{36}
\sol{B. A pie va $1 - \tfrac25 - \tfrac13 = \tfrac{4}{15}$: $\tfrac{4}{15}$ son 12, $\tfrac{1}{15}$ es 3 y el total 45.}
\item ¿Cuál fracción está entre $\tfrac13$ y $\tfrac12$?
\opciones{$\tfrac{3}{10}$}{$\tfrac{4}{7}$}{$\tfrac{5}{12}$}{$\tfrac{2}{9}$}
\sol{C. $\tfrac13 = \tfrac{4}{12} < \tfrac{5}{12} < \tfrac{6}{12} = \tfrac12$.}
\end{ejercicios}

\bitacora

% ------------------------------------------------------------------
\leccion{pu-cm-3}{Decimales y porcentajes}{Video «Porcentajes» y los juegos de las cien casillas y las cifras que se repiten.}

\capitulo{De cada 100}
«Por ciento» quiere decir \clave{de cada 100}. Una cuadrícula de 100 casillas es el total: pintar 35 es pintar el $35\,\%$.

\begin{casillas}\xcien{0}{35}{0}{$35\,\%$}\xcien{50}{100}{0}{$100\,\%$: el total}\xcien{100}{0}{0}{Tu turno: pinta el $65\,\%$}\end{casillas}

\turno{En la tercera cuadrícula pinta el $65\,\%$. ¿Qué porcentaje queda en blanco?}{El $35\,\%$: entre los dos completan el $100\,\%$.}

\capitulo{Tres formas de decir lo mismo}
Una fracción, un decimal y un porcentaje pueden nombrar la misma cantidad. Para pasar de una a otra se lleva a centésimas, o
se divide arriba entre abajo y se corre el punto dos lugares.

\figura{cien}{Veinticinco casillas de cien: la misma cantidad con tres nombres.}

\turno{Completa la tabla.\par\vspace{2pt}
{\renewcommand{\arraystretch}{1.5}\begin{tabularx}{\linewidth}{|>{\centering\arraybackslash}X|>{\centering\arraybackslash}X|>{\centering\arraybackslash}X|}\hline
\sffamily\small fracción & \sffamily\small decimal & \sffamily\small porcentaje\\\hline $\tfrac{7}{20}$ & & \\\hline & $0.08$ & \\\hline & & $62.5\,\%$\\\hline\end{tabularx}}}{$\tfrac{7}{20} = 0.35 = 35\,\%$;
$\tfrac{2}{25} = 0.08 = 8\,\%$; $\tfrac58 = 0.625 = 62.5\,\%$.}

\subsubsection{Decimales que no terminan}
Una fracción es una división: $\tfrac38 = 0.375$. Si el denominador, ya simplificado, solo tiene factores 2 y 5, el decimal
\clave{termina}; si tiene otro, es \clave{periódico}, como $\tfrac13 = 0.\overline{3}$ o $\tfrac17 = 0.\overline{142857}$. De vuelta:
si $x = 0.\overline{36}$, entonces $100x = 36.\overline{36}$; al restar, $99x = 36$ y $x = \tfrac{36}{99} = \tfrac{4}{11}$. El periodo va
sobre tantos nueves como cifras tenga.

\capitulo{El porcentaje de una cantidad}
Si la cuadrícula entera vale \$80\,000, cada casilla vale \$800. El $10\,\%$ son 10 casillas, \$8\,000: se corre el punto un
lugar. Con él se arman los demás: el $5\,\%$ es la mitad y $15\,\% = 10\,\% + 5\,\%$, es decir, \$12\,000.

\begin{casillas}\xcien{0}{10}{5}{$10\,\% + 5\,\%$}
\node[anchor=west,text width=92mm,align=left,font=\small,text=zink] at (44,-16) {De una vez: el porcentaje en decimal por la cantidad,
$0.15 \cdot 80\,000 = 12\,000$.\\[3pt] Al revés, ¿qué porcentaje es 18 de 72? La parte entre el todo: $18 \div 72 = 0.25$, el $25\,\%$.};\end{casillas}

Un truco más: el $p\,\%$ de $N$ es igual al $N\,\%$ de $p$. El $16\,\%$ de 25 es el $25\,\%$ de 16, que es 4.

\turno{La cuadrícula vale 360. ¿Cuánto es el $35\,\%$? ¿Y qué porcentaje es 54 de 360?}{Cada casilla vale 3.6:
$0.35 \cdot 360 = 126$. Y $54 \div 360 = 0.15$: el $15\,\%$.}

\capitulo{Subir y bajar}
Bajar un $30\,\%$ es quedarse con el $70\,\%$: multiplicar por $0.7$. Subir un $20\,\%$ es multiplicar por $1.2$. Los cambios
seguidos se \clave{multiplican}, porque el segundo se calcula sobre lo nuevo.

\begin{casillas}\xcien{0}{100}{0}{antes: $100\,\%$}\xcien{50}{70}{0}{baja $30\,\%$: queda $70\,\%$}\xcien{100}{70}{14}{sube $20\,\%$ de 70: $84\,\%$}\end{casillas}

Bajar 30 y subir 20 no es bajar 10: queda el $84\,\%$, es decir, bajó $16\,\%$. Para deshacer un cambio se divide: si \$56\,000 es
lo que quedó tras bajar un $30\,\%$, antes valía $56\,000 \div 0.7 = 80\,000$.

\turno{Un celular de \$900\,000 sube un $10\,\%$ y luego baja un $10\,\%$. ¿Vuelve a su precio?}{No: $900\,000 \cdot 1.1 \cdot 0.9 = 891\,000$;
bajó $1\,\%$.}

\begin{regla}
\begin{itemize}
  \item $p\,\% = \tfrac{p}{100}$. El $p\,\%$ de $N$ es $\tfrac{p}{100}\cdot N$.
  \item Fracciones que conviene saber: $\tfrac12 = 50\,\%$, $\tfrac14 = 25\,\%$, $\tfrac34 = 75\,\%$, $\tfrac15 = 20\,\%$,
  $\tfrac18 = 12.5\,\%$, $\tfrac13 \approx 33.3\,\%$, $\tfrac{1}{20} = 5\,\%$.
  \item Si un decimal se repite en bloques de $k$ cifras, la cifra del lugar $n$ es la misma que la del lugar
  $r$, donde $r$ es el residuo de dividir $n$ entre $k$ (si el residuo es 0, es la última del bloque).
\end{itemize}
\end{regla}

\subsection{Ejemplos resueltos}
\begin{ejemplo}[Comparar formatos]
Ordena de mayor a menor: $0.47$, $\tfrac{9}{20}$, $48\,\%$, $\tfrac{12}{25}$.
\solucion Todo a centésimas: $0.47$, $\tfrac{45}{100} = 0.45$, $0.48$, $\tfrac{48}{100} = 0.48$. Orden: $48\,\% =
\tfrac{12}{25} > 0.47 > \tfrac{9}{20}$. Dos de ellos son iguales, cosa que el examen aprovecha.
\end{ejemplo}

\begin{ejemplo}[El bloque que se repite]
¿Cuánto suman las primeras 40 cifras decimales de $\tfrac{1}{7}$?
\solucion $\tfrac17 = 0.\overline{142857}$, bloque de 6 cifras que suman $1+4+2+8+5+7 = 27$. Como $40 = 6 \cdot 6 + 4$,
hay seis bloques completos y las cuatro primeras cifras de otro: $6 \cdot 27 + (1+4+2+8) = 162 + 15 = 177$.
\end{ejemplo}

\begin{ejemplo}[Porcentaje inverso]
El $35\,\%$ de un número es 84. ¿Cuál es el número?
\solucion $35\,\%$ son 84, así que el $5\,\%$ es $84 \div 7 = 12$ y el $100\,\%$ es $12 \cdot 20 = 240$.
\end{ejemplo}

\begin{ejemplo}[¿Qué porcentaje es?]
En una encuesta, 27 de 180 personas prefieren el bus. ¿Qué porcentaje es?
\solucion $\tfrac{27}{180} = \tfrac{3}{20} = \tfrac{15}{100} = 15\,\%$.
\end{ejemplo}

\begin{trampa}
$0.3$ \textbf{no} es lo mismo que $0.\overline{3}$: el primero es $\tfrac{3}{10}$, el segundo $\tfrac13$. Y el
$0.5\,\%$ es $0.005$, no $0.5$. Lee los símbolos con lupa.
\end{trampa}

\begin{atajo}
Para porcentajes «feos», intercambia: el $18\,\%$ de 50 es el $50\,\%$ de 18, que es 9. Funciona siempre, porque
$\tfrac{18}{100}\cdot 50 = \tfrac{50}{100}\cdot 18$.
\end{atajo}

\begin{reto}
Al escribir $\tfrac37$ como decimal, ¿qué cifra queda en el lugar 100 después del punto?
\opciones{4}{2}{5}{7}
\end{reto}

\begin{retoresuelto}
$\tfrac37 = 0.\overline{428571}$: bloque de 6 cifras. Como $100 = 16 \cdot 6 + 4$, el lugar 100 es el cuarto del
bloque: \textbf{5}. \textbf{Opción C.} No hace falta escribir cien cifras; basta el residuo.
\end{retoresuelto}

\armar{
\arma{Los pedazos}{Arma el $37\,\%$ de 2\,000 usando solo el $50\,\%$, el $10\,\%$ y el $1\,\%$, sumando o restando.}{$50\,\% - 10\,\% - 3 \cdot 1\,\%
= 1\,000 - 200 - 60 = 740$. Error frecuente: sumar los porcentajes sin volverlos cantidades.}
\arma{La etiqueta vieja}{Inventa un precio de antes y un descuento para que la etiqueta nueva diga \$45\,000. Comprueba con la
multiplicación.}{Por ejemplo, \$60\,000 con $25\,\%$: $60\,000 \cdot 0.75 = 45\,000$. Error frecuente: sumarle el porcentaje al precio nuevo.}
\arma{La vuelta}{Un precio sube un $25\,\%$. ¿Qué descuento lo devuelve exactamente al inicio?}{$20\,\%$: hay que multiplicar por
$1 \div 1.25 = 0.8$. Error frecuente: responder $25\,\%$.}
}

\taller
\piso{Calienta}{conversiones}
\begin{ejercicios}[2]
\item $\tfrac{7}{8}$ en decimal \sol{$0.875$.}
\item $0.\overline{45}$ en fracción \sol{$\tfrac{45}{99} = \tfrac{5}{11}$.}
\item $12\,\%$ de $250$ \sol{$30$.}
\item $0.04$ en porcentaje \sol{$4\,\%$.}
\item $\tfrac{3}{40}$ en porcentaje \sol{$7.5\,\%$.}
\item $44\,\%$ de 25 \sol{$25\,\%$ de 44: $11$.}
\end{ejercicios}
\piso{Entrena}{problemas}
\begin{ejercicios}
\item ¿Cuál de estas fracciones da un decimal finito: $\tfrac{7}{12}$, $\tfrac{9}{40}$, $\tfrac{5}{6}$? \sol{$\tfrac{9}{40}$: $40 = 2^3 \cdot 5$. Da $0.225$.}
\item Un colegio tiene 640 estudiantes y el $15\,\%$ juega fútbol. ¿Cuántos no juegan fútbol? \sol{$85\,\%$ de 640: $544$.}
\item ¿Qué cifra ocupa el lugar 50 de $0.\overline{285714}$? \sol{$50 = 8\cdot 6 + 2$: la segunda cifra del bloque, 8.}
\item El 60\,\% de los asistentes a un concierto son 1\,500 personas. ¿Cuántos asistentes hay? \sol{$20\,\%$ son 500: total 2\,500.}
\end{ejercicios}
\piso{Como en el examen}{opción múltiple}
\begin{ejercicios}
\item ¿Cuál es el mayor?
\opciones{$\tfrac{5}{8}$}{$0.62$}{$63\,\%$}{$\tfrac{31}{50}$}
\sol{C. $0.625$, $0.62$, $0.63$, $0.62$.}
\item $0.\overline{9}$ es igual a
\opciones{$0.9$}{$1$}{$\tfrac{9}{10}$}{un número menor que $0.99$}
\sol{B. $\tfrac{9}{9} = 1$ (periodo sobre nueves).}
\item Si el $20\,\%$ de $A$ es igual al $30\,\%$ de $B$, entonces $A$ es
\opciones{el $150\,\%$ de $B$}{el $66\,\%$ de $B$}{el $110\,\%$ de $B$}{el $60\,\%$ de $B$}
\sol{A. $0.2A = 0.3B$, así que $A = 1.5B$.}
\end{ejercicios}

\bitacora

% ------------------------------------------------------------------
\leccion{pu-cm-4}{Potencias, raíces y estimación}{Video «Potencias» y los juegos de la torre de potencias y el muro de raíces.}

\capitulo{Multiplicar repetido}
Una potencia cuenta \clave{factores iguales}: $5^3 = 5 \cdot 5 \cdot 5 = 125$. No es $5 \cdot 3$. Cada factor más multiplica
todo otra vez, y con cero factores la potencia vale 1, el punto de partida de toda multiplicación.

\figura{potencias}{Una potencia no suma: duplica en cada paso.}

\turno{¿Cuántos factores 4 hacen falta para llegar a 256? Escribe el camino: $4,\ 16,\ \rule{9mm}{.4pt},\ \rule{9mm}{.4pt}$.}{Cuatro:
$4, 16, 64, 256$; $4^4 = 256$.}

\capitulo{El cuadrado como área}
Elevar \clave{al cuadrado} es armar un cuadrado: $4^2$ son 4 filas de 4 casillas. Si el lado se duplica, el área no se duplica:
se \emph{cuadruplica}.

\begin{casillas}\xcuad{0}{4}{4}{lado 4\\$4^2 = 16$}\xcuad{36}{8}{4}{lado 8\\$8^2 = 64$}
\node[anchor=west,text width=66mm,align=left,font=\small,text=zink] at (78,17) {En el de lado 8 caben cuatro cuadrados de lado 4.\\[3pt]
Con decimales es igual: si 100 casillas son 1, un lado de $0.4$ cubre 16 de 100: $0.4^2 = 0.16$.};\end{casillas}

\turno{Si el lado de un cuadrado se multiplica por 5, ¿por cuánto se multiplica su área? ¿Cuánto es $0.7^2$?}{Por 25. Y $0.7^2 = 0.49$.}

\capitulo{Contar factores}
Las leyes de las potencias no se memorizan: salen de contar. En $2^3 \cdot 2^4$ hay tres doses y cuatro doses: siete. En $(5^2)^3$
hay tres grupos de dos cincos: seis.

\begin{center}\renewcommand{\arraystretch}{1.35}
\begin{tabular}{@{}lll@{}}
\toprule
\textbf{Ley} & \textbf{Por qué} & \textbf{Ejemplo}\\\midrule
$a^m \cdot a^n = a^{m+n}$ & se juntan $m$ y $n$ factores & $2^3 \cdot 2^4 = 2^7$\\
$\dfrac{a^m}{a^n} = a^{m-n}$ & se cancelan $n$ factores & $\dfrac{3^7}{3^5} = 3^2$\\
$(a^m)^n = a^{mn}$ & $n$ grupos de $m$ factores & $(5^2)^3 = 5^6$\\
$(ab)^n = a^n b^n$ & se reparten los factores & $(2 \cdot 5)^4 = 2^4 \cdot 5^4$\\
$a^0 = 1,\ a^{-n} = \dfrac{1}{a^n}$ & seguir el patrón al dividir & $2^{-3} = \tfrac18$\\
\bottomrule
\end{tabular}
\end{center}

En las potencias de 10 se ve la cuenta: $10^n$ es un 1 con $n$ ceros y $10^n - 1$ son $n$ nueves. Como $2 \cdot 5 = 10$, cada
pareja de un 2 y un 5 produce un cero.

\turno{Escribe como una sola potencia: $\;6^4 \cdot 6^3$, $\;(3^5)^2\;$ y $\;\dfrac{7^9}{7^2}$.}{$6^7$, $3^{10}$ y $7^7$.}

\capitulo{La raíz, la pregunta al revés}
La \clave{raíz cuadrada} no pregunta cuánto da el área, sino qué lado tiene. $\sqrt{49} = 7$ porque un cuadrado de 49 casillas
tiene lado 7. Si no es exacta, se acorrala entre dos cuadrados: $64 < 72 < 81$, así que $8 < \sqrt{72} < 9$.

\begin{casillas}\xcuad{0}{8}{3}{$8^2 = 64$\\no alcanza}\xcuad{40}{9}{3}{$9^2 = 81$\\se pasa}
\node[anchor=west,text width=70mm,align=left,font=\small,text=zink] at (80,13) {Como $8.5^2 = 72.25$, $\sqrt{72}$ está apenas por debajo de 8.5.\\[3pt]
Para simplificar se saca el mayor cuadrado: $\sqrt{72} = \sqrt{36 \cdot 2} = 6\sqrt2$.};\end{casillas}

\textbf{Estimar.} Muchas preguntas no piden la cuenta exacta: $49.7 \cdot 20.3 \approx 50 \cdot 20 = 1\,000$. Si las opciones
están lejos unas de otras, estimar basta.

\turno{¿Entre qué enteros está $\sqrt{130}$ y de cuál está más cerca? Estima $\;6\,120 \cdot 0.49$.}{Entre 11 y 12
($121$ y $144$), más cerca de 11. El producto es casi la mitad de 6\,000: unos 3\,000.}

\begin{regla}
\begin{itemize}
  \item Si el lado de un cuadrado se multiplica por $k$, su área se multiplica por $k^2$; en un cubo, el volumen por $k^3$.
  \item $\sqrt{ab} = \sqrt a\,\sqrt b$, pero $\sqrt{a + b} \neq \sqrt a + \sqrt b$.
  \item $(-a)^n$ es positivo si $n$ es par y negativo si $n$ es impar.
  \item Cuadrados para tener a mano: $11^2 = 121,\ 12^2 = 144,\ 13^2 = 169,\ 14^2 = 196,\ 15^2 = 225,\ 16^2 = 256,
  \ 25^2 = 625$.
\end{itemize}
\end{regla}

\subsection{Ejemplos resueltos}
\begin{ejemplo}[Juntar en una base]
Simplifica $\dfrac{4^5 \cdot 8^2}{16^3}$.
\solucion Todo en base 2: $4^5 = 2^{10}$, $8^2 = 2^6$, $16^3 = 2^{12}$. Queda $2^{10 + 6 - 12} = 2^4 = 16$.
\end{ejemplo}

\begin{ejemplo}[El área crece al cuadrado]
El lado de un cuadrado aumenta un $30\,\%$. ¿En qué porcentaje aumenta su área?
\solucion El lado se multiplica por $1.3$ y el área por $1.3^2 = 1.69$: aumenta $69\,\%$, no $30\,\%$ ni $60\,\%$.
\end{ejemplo}

\begin{ejemplo}[Suma de cifras de una resta]
¿Cuánto suman las cifras de $10^{8} - 24$?
\solucion $10^8 - 24 = 99\,999\,976$: seis nueves, un 7 y un 6. La suma es $54 + 7 + 6 = 67$. Para verlo sin escribir:
$10^8 - 1$ son ocho nueves y luego se restan 23 a los dos últimos ($99 - 23 = 76$).
\end{ejemplo}

\begin{ejemplo}[Comparar potencias]
¿Cuál es mayor, $2^{30}$ o $3^{20}$?
\solucion Se lleva a un exponente común: $2^{30} = (2^3)^{10} = 8^{10}$ y $3^{20} = (3^2)^{10} = 9^{10}$. Como $9 > 8$,
es mayor $3^{20}$.
\end{ejemplo}

\begin{trampa}
$2^3$ no es $2 \cdot 3$, y $(a + b)^2$ no es $a^2 + b^2$: falta el término $2ab$. Por ejemplo, $(3 + 4)^2 = 49$, mientras
que $3^2 + 4^2 = 25$.
\end{trampa}

\begin{atajo}
Para contar cifras de $2^a \cdot 5^b$, empareja: el menor exponente da los ceros y lo que sobra se calcula aparte.
$2^7 \cdot 5^4 = 2^3 \cdot 10^4 = 80\,000$: cinco cifras.
\end{atajo}

\begin{reto}
¿Cuántas cifras tiene el número $2^{10} \cdot 5^{12}$?
\opciones{10}{11}{12}{22}
\end{reto}

\begin{retoresuelto}
Se emparejan diez 2 con diez 5: $2^{10} \cdot 5^{12} = (2 \cdot 5)^{10} \cdot 5^2 = 25 \cdot 10^{10}$. Es un 25 seguido
de diez ceros: $2 + 10 = 12$ cifras. \textbf{Opción C.} La D sale de sumar los exponentes, que no tiene sentido aquí.
\end{retoresuelto}

\armar{
\arma{Dos nombres}{Escribe $64$ como potencia de tres formas distintas.}{$2^6 = 4^3 = 8^2$ (también $64^1$). Error frecuente: $2 \cdot 32$, que
no es una potencia.}
\arma{Los exponentes justos}{Completa con enteros positivos: $\;3^{\square} \cdot 3^{\square} = 3^{8}\;$ y $\;(3^{\square})^{\square} = 3^{8}$, sin usar el 1.}{Sumas
de 8, como 3 y 5; y productos de 8: 2 y 4. Error frecuente: multiplicar los exponentes en el primer caso.}
\arma{El lado escondido}{Da un número que no sea cuadrado y cuya raíz esté entre 6 y 6.5. Explica cómo lo sabes.}{Cualquiera entre
36 y 42.25, como 40. Error frecuente: tomar la mitad en vez del lado.}
}

\taller
\piso{Calienta}{potencias y raíces}
\begin{ejercicios}[2]
\item $3^4$ \sol{$81$.}
\item $(-2)^5$ \sol{$-32$.}
\item $5^{-2}$ \sol{$\tfrac{1}{25}$.}
\item $\sqrt{98}$ simplificada \sol{$7\sqrt2$.}
\item $\dfrac{9^3}{27^2}$ \sol{$3^6 / 3^6 = 1$.}
\item $(2^3)^2 \cdot 2^{-4}$ \sol{$2^2 = 4$.}
\end{ejercicios}
\piso{Entrena}{problemas}
\begin{ejercicios}
\item ¿Entre qué dos enteros consecutivos está $\sqrt{150}$? \sol{Entre 12 y 13 ($144 < 150 < 169$).}
\item Si el radio de una pizza se duplica, ¿cuántas veces más área tiene? \sol{$2^2 = 4$ veces.}
\item ¿Cuántos ceros tiene al final $4^5 \cdot 25^3$? \sol{$2^{10} \cdot 5^6 = 2^4 \cdot 10^6$: seis ceros (es $16\,000\,000$).}
\item Estima $\dfrac{398 \cdot 51}{19.9}$. \sol{$\approx \tfrac{400 \cdot 50}{20} = 1\,000$.}
\end{ejercicios}
\piso{Como en el examen}{opción múltiple}
\begin{ejercicios}
\item $2^{20} + 2^{20}$ es igual a
\opciones{$2^{40}$}{$4^{20}$}{$2^{21}$}{$4^{40}$}
\sol{C. Dos veces $2^{20}$ es $2 \cdot 2^{20} = 2^{21}$.}
\item Si $3^x = 5$, entonces $9^{x+1}$ vale
\opciones{$45$}{$225$}{$75$}{$30$}
\sol{B. $9^{x+1} = 9 \cdot (3^x)^2 = 9 \cdot 25 = 225$.}
\item Un cubo de 3 cm de arista se agranda hasta 6 cm de arista. Su volumen se multiplica por
\opciones{2}{4}{6}{8}
\sol{D. $k = 2$ y el volumen se multiplica por $2^3 = 8$.}
\end{ejercicios}

\bitacora

% ------------------------------------------------------------------
\evaluacion
\begin{ejercicios}
\item El valor de $\;6 - 4 \div 2 \cdot (-3) + (-1)^{4}\;$ es
\opciones{$13$}{$-1$}{$7$}{$4$}
\sol{A. $4 \div 2 \cdot (-3) = -6$; $6 + 6 + 1 = 13$. (Lección 1.1)}
\item La temperatura en una ciudad varió así durante el día: $-3$, luego subió 8, bajó 11 y subió 4. ¿Qué temperatura marcaba al final?
\opciones{$-2$}{$2$}{$-4$}{$0$}
\sol{A. $-3 + 8 - 11 + 4 = -2$. (1.1)}
\item Ana se gasta $\tfrac13$ de su dinero y luego $\tfrac12$ de lo que le queda. ¿Qué fracción conserva?
\opciones{$\tfrac16$}{$\tfrac13$}{$\tfrac12$}{$\tfrac23$}
\sol{B. $\tfrac23 \cdot \tfrac12 = \tfrac13$. (1.2)}
\item ¿Cuál de las fracciones es la menor?
\opciones{$\tfrac{7}{9}$}{$\tfrac{4}{5}$}{$\tfrac{10}{13}$}{$\tfrac{3}{4}$}
\sol{D. A cada una le falta para 1: $\tfrac29, \tfrac15, \tfrac{3}{13}, \tfrac14$. A la menor le falta más: $\tfrac14$ es lo mayor. (1.2)}
\item ¿Cuántos medios vasos se llenan con $4\tfrac12$ litros si un vaso tiene $\tfrac34$ de litro?
\opciones{6}{9}{12}{18}
\sol{C. Medio vaso son $\tfrac38$ L; $\tfrac92 \div \tfrac38 = 12$. (1.2)}
\item El $15\,\%$ de 60 más el $60\,\%$ de 15 es
\opciones{$9$}{$18$}{$15$}{$4.5$}
\sol{B. Los dos valen 9. (1.3)}
\item La cifra 2026 después del punto en $0.\overline{123}$ es
\opciones{1}{2}{3}{0}
\sol{A. $2026 = 675 \cdot 3 + 1$: la primera del bloque. (1.3)}
\item En un curso de 45 estudiantes, 18 son hombres. ¿Qué porcentaje son mujeres?
\opciones{$40\,\%$}{$55\,\%$}{$60\,\%$}{$27\,\%$}
\sol{C. 27 de 45 es $\tfrac35 = 60\,\%$. (1.3)}
\item $\dfrac{2^{12} \cdot 3^{6}}{6^{6}}$ es igual a
\opciones{$2^6$}{$2^{18}$}{$6^{6}$}{$1$}
\sol{A. $6^6 = 2^6 3^6$; queda $2^{12-6} = 2^6$. (1.4)}
\item ¿Cuál número está más cerca de $\sqrt{90}$?
\opciones{$9$}{$9.5$}{$10$}{$8.5$}
\sol{B. $9.5^2 = 90.25$. (1.4)}
\item Si el lado de un cuadrado disminuye un $10\,\%$, su área disminuye
\opciones{$10\,\%$}{$20\,\%$}{$19\,\%$}{$81\,\%$}
\sol{C. $0.9^2 = 0.81$: queda el $81\,\%$, se pierde el $19\,\%$. (1.4)}
\item ¿Cuántas cifras tiene $8^{5} \cdot 5^{14}$?
\opciones{14}{15}{16}{19}
\sol{B. $2^{15} \cdot 5^{14} = 2 \cdot 10^{14}$: un 2 y catorce ceros, 15 cifras. (1.4)}
\end{ejercicios}
